\section{Výsledky experimentů}

\subsection{Cíl vyhodnocení}

Cílem této části je vyhodnotit chování čtyř navazujících simulací a ověřit hypotézy formulované v teoretické části. Důraz je kladen na srovnatelnost scénářů, proto jsou napříč simulacemi používány jednotné metriky a stejný postup vyhodnocení.

\subsection{Experimentální metodika}

\textbf{Princip experimentů:}
\begin{itemize}
\item pro každou simulaci je definován referenční scénář (výchozí nastavení)
\item následně je aplikován přístup \textit{one-factor-at-a-time} (OFAT)
\item pro vybrané dvojice parametrů jsou provedeny doplňkové kombinované scénáře
\item každá konfigurace je opakována pro více seedů
\end{itemize}

\textbf{Počet opakování:}
\begin{itemize}
\item každá konfigurace: \textbf{10} opakování (\texttt{seed = 1..10})
\end{itemize}

\textbf{Úrovně parametrů (low/mid/high):}
\begin{table}[h]
\centering
\begin{tabular}{l c c c}
\hline
\textbf{Parametr} & \textbf{Low} & \textbf{Mid} & \textbf{High} \\
\hline
\texttt{num-people} & 300 & 500 & 800 \\
\texttt{transmission-prob} & 30 & 60 & 85 \\
\texttt{incubation-days} & 2 & 5 & 9 \\
\texttt{mean-illness-duration} & 7 & 14 & 21 \\
\texttt{standard-speed} & 1 & 2 & 3 \\
\texttt{fatality-rate} & 0.5 & 1.8 & 4.0 \\
\texttt{healthcare-capacity} & 80 & 140 & 240 \\
\texttt{overload-multiplier} & 1.2 & 2.0 & 3.0 \\
\texttt{resource-stock} & 600 & 1400 & 2600 \\
\texttt{resource-regeneration} & 2 & 8 & 20 \\
\texttt{treatment-cost} & 6 & 12 & 24 \\
\texttt{resource-penalty} & 1.2 & 1.8 & 2.8 \\
\texttt{vaccination-rate} & 20 & 50 & 80 \\
\texttt{high-risk-share} & 10 & 20 & 35 \\
\hline
\end{tabular}
\caption{Doporučené úrovně parametrů pro OFAT a kombinované scénáře}
\label{tab:param_levels}
\end{table}

\textbf{Délka běhu:}
\begin{itemize}
\item simulace běží do stabilizace nebo do pevného limitu 8760 ticků (dle implementace modelu)
\end{itemize}

\subsection{Sledované metriky}

V každém scénáři jsou sledovány následující metriky:
\begin{itemize}
\item \textbf{Peak infected} -- maximální počet současně infekčních agentů
\item \textbf{Total deaths} -- celkový počet úmrtí
\item \textbf{Time to stabilization} -- čas do stabilizace systému
\item \textbf{Wave count / reinfections} -- počet vln nebo indikátor opětovných nákaz
\end{itemize}

U modelů se zdroji se navíc sleduje:
\begin{itemize}
\item \textbf{Resource depletion ratio} -- podíl času s nízkou nebo nulovou zásobou zdrojů
\item \textbf{Over-capacity load} -- přetížení kapacity systému v čase
\end{itemize}

\subsection{Referenční scénáře}

Tabulka~\ref{tab:baseline} uvádí výchozí konfigurace použité jako baseline.

\begin{table}[h]
\centering
\begin{tabular}{l l}
\hline
\textbf{Model} & \textbf{Referenční konfigurace} \\
\hline
Simulace 1 & výchozí hodnoty sliderů \\
Simulace 2 & výchozí hodnoty sliderů \\
Simulace 3 & výchozí hodnoty sliderů \\
Simulace 4 & výchozí hodnoty sliderů \\
\hline
\end{tabular}
\caption{Referenční scénáře (baseline)}
\label{tab:baseline}
\end{table}

\subsection{Minimální sada scénářů}

Aby byl experiment zvládnutelný, ale zároveň obhajitelný, je doporučena následující minimální sada:
\begin{itemize}
\item baseline pro každý model
\item OFAT: 4 klíčové parametry na model, každý ve 3 úrovních (low/mid/high)
\item 2 kombinované scénáře na model (interakce parametrů)
\item 10 seedů na každou konfiguraci
\end{itemize}

Tím vznikne přibližně 15 konfigurací na model, tj. cca 150 běhů na model a 600 běhů celkem.

\subsection{Výsledky Simulace 1}

\subsubsection{Baseline výsledek (\texttt{baseline\_auto}, 10 běhů)}

\begin{table}[h]
\centering
\begin{tabular}{l c}
\hline
\textbf{Metrika} & \textbf{Hodnota (mean)} \\
\hline
Final susceptible & 0.0 \\
Final exposed & 0.0 \\
Final infected & 0.0 \\
Final recovered & 500.0 \\
Ever infected & 500.0 \\
Deaths & 0.0 \\
Time to stabilization (ticks) & 1327.4 \\
\hline
\end{tabular}
\caption{Baseline výsledky Simulace 1}
\label{tab:s1_baseline}
\end{table}

\subsubsection{OFAT experimenty}

V Simulaci 1 byly měněny parametry:
\begin{itemize}
\item \texttt{transmission-prob}
\item \texttt{incubation-days}
\item \texttt{standard-speed}
\item \texttt{num-people}
\end{itemize}

\begin{table}[h]
\centering
\begin{tabular}{l c c c c}
\hline
\textbf{Parametr} & \textbf{Úroveň} & \textbf{Peak infected} & \textbf{Time to stab.} & \textbf{Poznámka} \\
\hline
\texttt{transmission-prob} & low (30) & 0.0 (final I) & 1596.8 & nejpomalejší do stabilizace \\
\texttt{transmission-prob} & mid (60) & 0.0 (final I) & 1369.5 & referenční průběh \\
\texttt{transmission-prob} & high (85) & 0.0 (final I) & 1285.7 & nejrychlejší do stabilizace \\
\hline
\end{tabular}
\caption{Ukázka OFAT výsledků pro Simulaci 1}
\label{tab:s1_ofat}
\end{table}

\textbf{Interpretace:}  
Baseline konfigurace vede k úplnému promoření populace bez úmrtí; systém se stabilizuje přibližně po 1327 tickech. OFAT potvrzuje, že s vyšším \texttt{transmission-prob} se epidemie uzavírá rychleji (kratší doba do stabilizace).

\subsection{Výsledky Simulace 2}

\subsubsection{Baseline výsledek (\texttt{baseline\_auto}, 10 běhů)}

\begin{table}[h]
\centering
\begin{tabular}{l c}
\hline
\textbf{Metrika} & \textbf{Hodnota (mean)} \\
\hline
Final susceptible & 0.0 \\
Final exposed & 0.0 \\
Final infected & 0.0 \\
Final recovered & 484.1 \\
Ever infected & 500.0 \\
Deaths & 15.9 \\
Time to stabilization (ticks) & 1332.2 \\
\hline
\end{tabular}
\caption{Baseline výsledky Simulace 2}
\label{tab:s2_baseline}
\end{table}

\subsubsection{Vliv kapacity systému}

Zde je klíčový vztah mezi \texttt{healthcare-capacity}, \texttt{overload-multiplier} a \texttt{fatality-rate}.

\begin{table}[h]
\centering
\begin{tabular}{c c c c}
\hline
\textbf{Capacity} & \textbf{Overload mult.} & \textbf{Total deaths} & \textbf{Over-capacity load} \\
\hline
low (80) & low (1.2) & 11.8 & N/A \\
low (80) & high (3.0) & 26.6 & N/A \\
high (240) & low (1.2) & 10.7 & N/A \\
high (240) & high (3.0) & 16.9 & N/A \\
\hline
\end{tabular}
\caption{Citlivost Simulace 2 na kapacitu a přetížení}
\label{tab:s2_capacity}
\end{table}

\textbf{Interpretace:}  
Při výchozím nastavení vzniká nenulová úmrtnost (průměrně 15.9 úmrtí), zatímco v Simulaci 1 byla nulová. Grid scénář potvrzuje, že nejhorší kombinace je nízká kapacita + vysoký overload multiplikátor.

\subsection{Výsledky Simulace 3}

\subsubsection{Baseline výsledek (\texttt{baseline\_auto}, 10 běhů)}

\begin{table}[h]
\centering
\begin{tabular}{l c}
\hline
\textbf{Metrika} & \textbf{Hodnota (mean)} \\
\hline
Final susceptible & 0.0 \\
Final exposed & 42.8 \\
Final infected & 124.6 \\
Final recovered & 218.6 \\
Ever infected & 500.0 \\
Deaths & 114.0 \\
Final resources & 1396.8 \\
Runtime (ticks) & 8760.0 \\
\hline
\end{tabular}
\caption{Baseline výsledky Simulace 3}
\label{tab:s3_baseline}
\end{table}

\subsubsection{Vliv omezených zdrojů}

Byly testovány parametry \texttt{resource-stock}, \texttt{resource-regeneration}, \texttt{treatment-cost} a \texttt{resource-penalty}.

\begin{table}[h]
\centering
\begin{tabular}{c c c c c}
\hline
\textbf{Stock} & \textbf{Regen} & \textbf{Cost} & \textbf{Deaths} & \textbf{Resource depletion} \\
\hline
low (600) & low (2) & high (24) & 145.3 & vysoká (final res. 17.6) \\
low (600) & high (20) & high (24) & 109.0 & nízká až střední (final res. 574.8) \\
high (2600) & low (2) & low (6) & 114.0 & střední (final res. 15.6) \\
high (2600) & high (20) & low (6) & 109.0 & velmi nízká (final res. 2599.4) \\
\hline
\end{tabular}
\caption{Vliv zdrojových parametrů v Simulaci 3}
\label{tab:s3_resources}
\end{table}

\textbf{Interpretace:}  
Model v baseline nedosáhne stabilizace před limitem 8760 ticků a vykazuje perzistentní infekci (nenulové E i I na konci běhu), což je konzistentní s chováním modelu s reinfekcí a omezenými zdroji.

\subsection{Výsledky Simulace 4}

\subsubsection{Baseline výsledek (\texttt{baseline\_auto}, 10 běhů)}

\begin{table}[h]
\centering
\begin{tabular}{l c}
\hline
\textbf{Metrika} & \textbf{Hodnota (mean)} \\
\hline
Final susceptible & 25.0 \\
Final exposed & 33.1 \\
Final infected & 98.2 \\
Final recovered & 284.6 \\
Ever infected & 475.0 \\
Deaths & 34.2 \\
Final resources & 1394.8 \\
Runtime (ticks) & 8760.0 \\
\hline
\end{tabular}
\caption{Baseline výsledky Simulace 4}
\label{tab:s4_baseline}
\end{table}

\subsubsection{Porovnání strategií}

Strategie jsou reprezentovány parametrem \texttt{strategy-mode}. Dále byl sledován vliv \texttt{high-risk-share} a \texttt{vaccination-rate}.

\begin{table}[h]
\centering
\begin{tabular}{c c c c c}
\hline
\textbf{Strategy mode} & \textbf{High-risk share} & \textbf{Vacc. rate} & \textbf{Deaths} & \textbf{Peak infected} \\
\hline
0 & low (10) & low (20) & 74.7 & 124.1 (final I) \\
1 & low (10) & low (20) & 103.6 & 100.5 (final I) \\
2 & low (10) & low (20) & 76.4 & 119.0 (final I) \\
0 & high (35) & high (80) & 14.5 & 80.1 (final I) \\
1 & high (35) & high (80) & 21.4 & 78.6 (final I) \\
2 & high (35) & high (80) & 14.3 & 78.6 (final I) \\
\hline
\end{tabular}
\caption{Porovnání strategií alokace v Simulaci 4}
\label{tab:s4_strategy}
\end{table}

\textbf{Interpretace:}  
V baseline konfiguraci je celkový podíl nakažených nižší než v Simulaci 3 (ever infected 475 vs. 500), současně výrazně klesá i úmrtnost (34.2 vs. 114.0), což naznačuje přínos kombinace očkování a strategie alokace zdrojů.

\subsection{Srovnání modelů}

\begin{table}[h]
\centering
\begin{tabular}{l c c c c}
\hline
\textbf{Model} & \textbf{Peak infected} & \textbf{Deaths} & \textbf{Time to stab.} & \textbf{Poznámka} \\
\hline
Simulace 1 & 0.0 & 0.0 & 1327.4 & baseline \\
Simulace 2 & 0.0 & 15.9 & 1332.2 & kapacita systému \\
Simulace 3 & 124.6 (final I) & 114.0 & 8760.0 (limit) & zdroje + reinfekce \\
Simulace 4 & 98.2 (final I) & 34.2 & 8760.0 (limit) & strategie + očkování \\
\hline
\end{tabular}
\caption{Souhrnné porovnání všech modelů}
\label{tab:all_models}
\end{table}

\subsection{Vyhodnocení hypotéz}

\textbf{H1:} Zvýšení hustoty populace vede k nelineárnímu růstu rychlosti šíření.  
Průběžně: \textbf{čeká na OFAT scénáře} (zatím vyhodnocen pouze baseline).

\textbf{H2:} Omezené zdroje zpomalují šíření, ale mohou způsobit lokální kolapsy.  
Průběžně: \textbf{částečně potvrzeno} baseline porovnáním Simulace 2 vs. 3 (vyšší perzistence infekce a výrazný nárůst úmrtnosti při zdrojové dynamice).

\textbf{H3:} Heterogenita agentů stabilizuje systém.  
Průběžně: \textbf{částečně potvrzeno} baseline porovnáním Simulace 3 vs. 4 (nižší úmrtnost a nižší celková infekce při heterogenních strategiích).

\subsection{Shrnutí}

TODO: 1--2 odstavce se závěrečnou interpretací:
\begin{itemize}
\item hlavní zjištění
\item praktický význam výsledků
\item omezení experimentu
\item návrhy pro další rozšíření
\end{itemize}

\subsection{Automatizace a replikovatelnost}

Pro plně replikovatelný běh je doporučeno využít \textbf{BehaviorSpace} v headless režimu:
\begin{itemize}
\item definovat experimenty v NetLogo (pojmenované scénáře)
\item spouštět je skriptem s fixním počtem opakování a seedů
\item exportovat výsledky do CSV
\item agregovat průměr, směrodatnou odchylku a percentily
\end{itemize}

Praktický postup a skripty jsou uvedeny v souborech:
\begin{itemize}
\item \texttt{automation.md}
\item \texttt{run\_behaviorspace.ps1}
\item \texttt{aggregate\_results.py}
\end{itemize}
